\section{领域四：Prompt 工程与结构化输出（20\%）}

这个领域的核心原则只有两个字：\textbf{明确}（be explicit）。

\subsection{模糊指令 vs 明确指令}

\begin{warningbox}{模糊指令无法提升精度}
以下常见指令实际上不起作用：
\begin{itemize}[leftmargin=1.5em]
    \item "Be conservative"——不会真正提升精度
    \item "Only report high-confidence findings"——不会真正减少误报
\end{itemize}
\textbf{有效做法}：精确定义哪些问题应该上报、哪些应该跳过，并为每个严重级别提供具体的代码示例。
\end{warningbox}

\subsection{Few-shot 示例的杠杆效应}

\begin{importantbox}{Few-shot 是考试中测试的最高杠杆技术}
使用 2--4 个精心设计的示例，展示以下内容：
\begin{itemize}[leftmargin=1.5em]
    \item 模糊情况的处理方式
    \item 为什么选择了某个动作而非其他替代方案
    \item 推理过程的展示
\end{itemize}
关键不在于示例数量，而在于选择\textbf{模糊边界案例}（ambiguous cases）作为示例。
\end{importantbox}

\subsection{tool\_use 与 JSON Schema}

\texttt{tool\_use} 配合 JSON schema 可以消除\textbf{语法错误}，但\textbf{无法消除语义错误}。在 schema 设计中需要注意：

\begin{itemize}[leftmargin=1.5em]
    \item \textbf{nullable 字段}：当源数据可能缺失时，使用 nullable 字段而非让模型编造值
    \item \textbf{"unclear" 枚举值}：为模型提供一个"不确定"的选项
    \item \textbf{"other" + detail 字符串}：覆盖预定义枚举之外的情况
\end{itemize}

\subsection{Message Batches API}

\begin{knowledgebox}{Batches API 的特性与适用场景}
\begin{itemize}[leftmargin=1.5em]
    \item 节省 50\% 成本
    \item 处理时间最长 24 小时，无延迟 SLA
    \item 不支持多轮 tool calling
    \item \textbf{适用}：隔夜批量报告生成
    \item \textbf{不适用}：需要阻塞等待结果的 pre-merge 检查（应使用同步 API）
\end{itemize}
\end{knowledgebox}

\subsection{推荐实践项目}

\begin{knowledgebox}{动手练习：结构化数据抽取管道}
构建一个使用 \texttt{tool\_use} 的数据抽取管道：
\begin{itemize}[leftmargin=1.5em]
    \item 定义包含 required、optional 和 nullable 字段的 JSON schema
    \item 实现 validation-retry 循环
    \item 通过 Batches API 批量处理
    \item 按 \texttt{custom\_id} 处理失败情况
\end{itemize}
\end{knowledgebox}

\subsection{本章小结}

Prompt 工程的核心是"明确"：用具体标准替代模糊指令，用 2--4 个边界案例的 few-shot 示例来引导模型行为。结构化输出方面，\texttt{tool\_use} + JSON schema 解决语法问题，nullable 字段和 "unclear" 枚举值解决语义问题。Batches API 适合非实时批量任务。

